% =============================================================================
%  BlueCCGL — Project Plan (Crisis Framing)
%  Bluesky Contrastive Cascade Graph Learning
%  Applied to the Israel-Palestine Conflict
% =============================================================================
\documentclass[12pt,a4paper]{article}

% ---- Packages ---------------------------------------------------------------
\usepackage[utf8]{inputenc}
\usepackage[T1]{fontenc}
\usepackage{lmodern}
\usepackage[english]{babel}
\usepackage{geometry}
\geometry{margin=2.5cm}

\usepackage{amsmath,amssymb}
\usepackage{graphicx}
\usepackage{booktabs}          % professional tables
\usepackage{longtable}         % tables spanning pages
\usepackage{array}             % extended column definitions
\usepackage{enumitem}          % list customisation
\usepackage{xcolor}            % colour for warning flags
\usepackage{hyperref}
\hypersetup{
  colorlinks  = true,
  linkcolor   = blue!70!black,
  citecolor   = green!50!black,
  urlcolor    = blue!60!black,
}
\usepackage[
  backend=biber,
  style=apa,
  sorting=nyt,
  natbib=true
]{biblatex}
\addbibresource{references.bib}

% ---- Title ------------------------------------------------------------------
\title{\textbf{BlueCCGL: Unsupervised Discovery of Crisis Information Cascade
       Archetypes on Bluesky}\\[6pt]
       \large The Israel-Palestine Conflict as a Natural Experiment\\
       in Decentralised Crisis Communication\\[6pt]
       \large Project Plan}
\author{}
\date{}

% =============================================================================
\begin{document}
\maketitle

% =============================================================================
\section{Introduction}
% =============================================================================

Understanding how crisis information propagates through social networks has
become a central question in computational social science.  During acute
geopolitical crises, information cascades---the chains of reposts through which
content reaches new audiences---exhibit dynamics that differ fundamentally from
steady-state discourse: urgency compresses decision-making, emotional intensity
amplifies sharing, and community boundaries become more or less permeable
depending on the nature of the event.  On centralised platforms such as
Twitter/X and Weibo, crisis cascade dynamics have been studied extensively
using supervised deep learning methods
\citep{chen2019cascn, cao2017deephawkes, yuan2024casflow}.  However, these
approaches share two critical limitations: they require ground-truth labels
(e.g., ``viral'' vs.\ ``non-viral''), and they are designed for platforms where
a single opaque recommendation algorithm governs content visibility.  Neither
assumption holds on decentralised social networks, and neither is compatible
with the inherent unpredictability of crisis information dynamics.

Bluesky, built on the AT Protocol \citep{kleppmann2024atprotocol}, represents a
new class of social platform where users can choose or disable algorithmic feeds
entirely, where identity is portable across servers, and where the full firehose
of platform activity is publicly accessible.  During the period covered by the
\citet{failla2024bluesky} Zenodo dataset (February~2023 -- March~2024), the
Israel-Palestine conflict escalated dramatically following the October~7, 2023
Hamas attack, producing sustained, multi-phase crisis discourse.  Bluesky's
architecture creates a unique natural experiment: crisis content spread both
organically through follower networks and via user-selected algorithmic feeds
(notably the News feed, which spiked in engagement during October~2023)---and
crucially, the structural trace of both pathways is observable.  The platform's
user base during this period was predominantly English-speaking, Western, and
politically left-leaning \citep{quelle2025bluesky}, making it a study of how an
ideologically aligned early-adopter community processes crisis information
through decentralised infrastructure.  Yet no study has analysed crisis
information cascade dynamics on Bluesky or any AT~Protocol platform at a
structural level.

This research proposes \textbf{BlueCCGL} (Bluesky Contrastive Cascade Graph
Learning), a framework for the unsupervised discovery of crisis information
cascade archetypes on Bluesky, applied to the Israel-Palestine conflict
(October~7, 2023 -- March~18, 2024).  BlueCCGL extends the Contrastive Cascade
Graph Learning (CCGL) framework of \citet{xu2022ccgl} with four novel
contributions:
\begin{enumerate}[label=\textbf{(\arabic*)}]
  \item an upstream \emph{NLP filtering pipeline} that identifies
        conflict-related posts via multilingual keyword matching and isolates
        cascades across six major escalation events;
  \item \emph{feed-aware graph augmentation strategies} that respect the
        semantic distinction between organic and algorithmically-amplified
        spread;
  \item \emph{edge-type aware message passing} via Relational Graph
        Convolutional Networks \citep[R-GCN;][]{schlichtkrull2018rgcn} that
        explicitly encodes whether a cascade edge represents follower-based or
        non-follower-based propagation; and
  \item \emph{hierarchical graph pooling} \citep{ying2018diffpool} enriched
        with sentiment composition features that capture multi-scale cascade
        structure from individual user decisions to community-level patterns.
\end{enumerate}
By learning cascade representations through contrastive self-supervised learning
and clustering them with HDBSCAN \citep{mcinnes2017hdbscan}, BlueCCGL discovers
natural crisis cascade archetypes without requiring any labelled data.

% =============================================================================
\section{Related Work}
% =============================================================================

% -----------------------------------------------------------------------------
\subsection{Information Cascade Prediction}
% -----------------------------------------------------------------------------

Graph neural network approaches to cascade modelling have evolved rapidly.
\textbf{CasCN} \citep{chen2019cascn} was among the first to apply GNNs to
cascade prediction, modelling cascades as sequences of graph snapshots with
temporal attention.  \textbf{DeepHawkes} \citep{cao2017deephawkes} bridged
classical Hawkes processes with deep learning for interpretable attention over
user influence.  \textbf{CasFlow} \citep{yuan2024casflow} introduced a
variational approach with hierarchical pooling, achieving state-of-the-art
performance on Weibo and Twitter.  \textbf{HierCas} \citep{zhang2024hiercas}
extended this with multi-resolution modelling capturing both user-level and
community-level dynamics.

\paragraph{Gap.}
All of the above methods are \emph{supervised}, requiring ground-truth
popularity or virality labels.  None account for heterogeneous edge types
(e.g., follower vs.\ non-follower propagation), none have been applied to
crisis-specific cascades, and all were evaluated exclusively on centralised
platforms.  BlueCCGL addresses all three gaps: it is unsupervised, edge-type
aware, crisis-focused, and targets a decentralised platform.

% -----------------------------------------------------------------------------
\subsection{Contrastive and Self-Supervised Graph Learning}
% -----------------------------------------------------------------------------

\textbf{CCGL} \citep{xu2022ccgl} is the direct methodological predecessor of
this work.  It introduced contrastive self-supervised pre-training for cascade
graphs using sub-cascade sampling for augmentation and an InfoNCE graph-level
contrastive loss, demonstrating significant improvement over supervised
baselines.  \textbf{GraphCL} \citep{you2020graphcl} provided a systematic study
of graph augmentation strategies (node dropping, edge perturbation, attribute
masking, subgraph sampling) that informs augmentation design.  \textbf{GRACE}
\citep{zhu2020grace} introduced node-level corruption-based contrastive
learning.

\paragraph{Gap.}
CCGL's augmentations are generic (random sub-cascade sampling) and carry no
domain semantics.  No existing contrastive graph learning method incorporates
augmentations that are aware of the distinction between organic and
algorithmically-mediated spread, nor has any been specifically designed for
crisis information cascades.  BlueCCGL introduces feed-aware augmentation
strategies (follower-subgraph extraction, temporal prefix sampling, bridge-node
removal) designed to preserve the semantic structure of crisis-driven
propagation.

% -----------------------------------------------------------------------------
\subsection{Heterogeneous GNNs and Organic vs.\ Algorithmic Spread}
% -----------------------------------------------------------------------------

\textbf{R-GCN} \citep{schlichtkrull2018rgcn} introduced relation-specific
weight matrices with basis decomposition for parameter efficiency.  \textbf{HAN}
\citep{wang2019han} proposed hierarchical attention across node-level and
semantic-level, while \textbf{HGT} \citep{hu2020hgt} applied Transformer-style
attention to heterogeneous graphs.  \textbf{RHGNN} \citep{yu2023rhgnn}
demonstrated the importance of modelling multiple relationship types in social
network analysis.

On the empirical side, \citet{huszar2022twitter} provided a landmark internal
Twitter study showing algorithmic amplification of political content---a gold
standard requiring platform-insider access.  \citet{goel2016structural} defined
the ``structural virality'' metric (Wiener index) to distinguish broadcast from
viral spread, and \citet{vosoughi2018spread} showed that true and false news
produce structurally different cascades.  \citet{bandy2021auditing} developed
methodology for auditing algorithmic curation.

\paragraph{Gap.}
No prior work has applied heterogeneous GNN encoding to crisis information
cascades, nor has any study examined organic vs.\ algorithmic spread of
conflict-related content on a decentralised platform.  BlueCCGL exploits
Bluesky's observable follow graph to directly label cascade edges and uses R-GCN
encoding to learn relation-aware representations of crisis cascades.

% -----------------------------------------------------------------------------
\subsection{Decentralised Social Networks and Bluesky}
% -----------------------------------------------------------------------------

Research on decentralised social networks has focused primarily on the
Fediverse.  \citet{zignani2018mastodon} conducted the first large-scale Mastodon
measurement study, and \citet{raman2019challenges} analysed technical challenges
of federation protocols.  For the AT~Protocol, \citet{kleppmann2024atprotocol}
described Bluesky's architecture including its Firehose, and
\citet{quelle2025bluesky} showed that Bluesky users share predominantly
left-centre news sources and exhibit significant issue-based polarisation on the
Israel-Palestine conflict.

Most notably, \citet{failla2024bluesky} released the first large-scale Bluesky
dataset, covering over 4M~users ($\sim$81\% of registered accounts), 235M~posts,
a directed follower graph of $\sim$144M~edges, and timestamped interaction
graphs---all on Zenodo
(DOI:~\href{https://doi.org/10.5281/zenodo.14669616}{10.5281/zenodo.14669616}).
The dataset includes posts from 11~feed generators and user feed-subscription
data, providing ground-truth information about algorithmic exposure.  The paper
notes a spike in News feed engagement during October~2023, attributable to the
Israel-Palestine conflict.

\paragraph{Gap.}
While \citet{failla2024bluesky} provide the first comprehensive Bluesky dataset
with descriptive analysis, no study has performed cascade-level \emph{structural}
analysis on Bluesky or any AT~Protocol platform---and no study has examined
crisis-specific cascade dynamics on a decentralised platform.  BlueCCGL fills
both gaps.

% -----------------------------------------------------------------------------
\subsection{Hierarchical Graph Pooling and Clustering}
% -----------------------------------------------------------------------------

\textbf{DiffPool} \citep{ying2018diffpool} introduced differentiable soft
cluster assignment for hierarchical graph representation learning.
\textbf{SAGPool} \citep{lee2019sagpool} provided a more efficient
attention-based alternative, and \textbf{MinCutPool}
\citep{bianchi2020mincutpool} offered a spectral-clustering-grounded approach.
For unsupervised pattern discovery, \textbf{HDBSCAN} \citep{mcinnes2017hdbscan}
provides density-based clustering without pre-specifying cluster count, while
\textbf{DEC} \citep{xie2016dec} and \textbf{DMoN} \citep{tsitsulin2023dmon}
offer end-to-end deep clustering alternatives.

\paragraph{Gap.}
Hierarchical pooling has not been combined with contrastive learning for cascade
graphs, nor applied with sentiment-enriched node features in a crisis context.
BlueCCGL integrates DiffPool-style hierarchical pooling with sentiment
composition as a graph-level feature, enabling multi-scale representations that
capture how crisis sentiment distributes across cascade layers.

% =============================================================================
\section{Research Questions}
% =============================================================================

\paragraph{Main Research Question.}
\begin{quote}
\emph{Can unsupervised contrastive graph learning discover meaningful crisis
information cascade archetypes on Bluesky that reflect the interplay between
organic follower-based spread and algorithmic amplification during the
Israel-Palestine conflict?}
\end{quote}

\noindent
This is decomposed into \textbf{socially-grounded} questions (what do we learn
about crisis communication?) and \textbf{methodological} questions (do the
technical extensions contribute?).

\paragraph{SRQ1: Crisis Phase Evolution.}
Do crisis cascade archetypes shift across conflict phases---e.g., from
broadcast-dominated cascades during acute shock (Oct~7) to deeper organic-viral
cascades during sustained advocacy (Feb--Mar~2024)?

\paragraph{SRQ2: Event-Type Structural Signatures.}
Do different crisis event categories (military escalation, humanitarian crisis,
legal milestone, individual story, ceasefire) produce structurally distinct
cascade archetypes?

\paragraph{SRQ3: Algorithmic Amplification.}
What proportion of crisis cascade edges are follower-based (organic) versus
non-follower-based (suggesting algorithmic or cross-community exposure), and does
this ratio differ across mega-events and archetypes?

\paragraph{SRQ4: Sentiment-Structure Coupling.}
Is the sentiment composition of a crisis cascade correlated with its structural
archetype?

\paragraph{MRQ1: Archetype Discovery.}
What natural cascade archetypes emerge from HDBSCAN clustering of contrastive
cascade embeddings, and how do they differ in structural properties (depth,
breadth, branching factor, follower-edge ratio, sentiment composition)?

\paragraph{MRQ2: Edge-Type Aware Encoding.}
How does R-GCN encoding of follower vs.\ non-follower edges affect
representation quality compared to homogeneous GNN baselines?

\paragraph{MRQ3: Feed-Aware Augmentation.}
Do domain-specific augmentation strategies produce higher-quality representations
than generic graph augmentations, as measured by clustering coherence and
downstream task performance?

% =============================================================================
\section{Methodology}
% =============================================================================

\subsection{NLP Filtering Pipeline}

The filtering pipeline identifies conflict-related posts from the full
235M-post corpus and assigns them to mega-event temporal windows.  Keywords are
organised into five tiers: (1)~high-precision English core terms (e.g.,
\emph{gaza}, \emph{hamas}, \emph{ceasefire}); (2)~English hashtag patterns
(e.g., \emph{\#FreePalestine}, \emph{\#CeasefireNow}); (3)~Arabic keywords;
(4)~Hebrew keywords; and (5)~supplementary named entities.  A post is tagged as
conflict-related if it matches $\geq$1~keyword from any tier, applied via
compiled regex on the \texttt{text} field.

Only posts dated October~1, 2023 -- March~18, 2024 are considered.  Each
conflict-related post is assigned to one of six mega-event windows
(Table~\ref{tab:events}) if its timestamp falls within a $\pm$3-day range, or
labelled \textbf{baseline} otherwise.

\begin{longtable}{
  @{}
  >{\centering\arraybackslash}p{0.5cm}
  >{\raggedright\arraybackslash}p{3.8cm}
  >{\centering\arraybackslash}p{2.5cm}
  >{\raggedright\arraybackslash}p{3.0cm}
  >{\raggedright\arraybackslash}p{4.0cm}
  @{}
}
\caption{Mega-events used as temporal anchors for cascade comparison.}
\label{tab:events} \\
\toprule
\textbf{\#} & \textbf{Event} & \textbf{Date(s)} & \textbf{Category} & \textbf{Rationale} \\
\midrule
\endfirsthead
\toprule
\textbf{\#} & \textbf{Event} & \textbf{Date(s)} & \textbf{Category} & \textbf{Rationale} \\
\midrule
\endhead
\bottomrule
\endlastfoot

E1 & Hamas attack / Israel declares war & Oct 7, 2023 & Military (acute shock) &
  Founding event; maximum uncertainty \\
\addlinespace
E2 & Al-Ahli hospital explosion & Oct 17, 2023 & Humanitarian / contested &
  Biggest Oct.\ flashpoint; contested narratives \\
\addlinespace
E3 & Qatar-brokered ceasefire \& hostage exchanges & Nov 24--30, 2023 &
  Diplomatic & Multi-day sustained event; hope-inflected discourse \\
\addlinespace
E4 & ICJ genocide ruling & Jan 26, 2024 & Legal / institutional &
  Legal/moral framing; deliberative cascades expected \\
\addlinespace
E5 & Aaron Bushnell self-immolation & Feb 25, 2024 & Individual story &
  Highly viral; emotionally resonant \\
\addlinespace
E6 & Flour massacre + 30{,}000 death toll & Feb 29, 2024 & Humanitarian milestone &
  Compound event; post-public-launch user base \\

\end{longtable}

\subsection{Cascade Construction}

Cascade roots are original (non-repost) conflict-related posts.  Downstream
reposts are traced via \texttt{repost\_from}/\texttt{reposted\_author} fields
and \texttt{interactions.csv.gz}.  For each reposter, a parent is assigned by
priority: (1)~a followed user who already reposted earlier (organic path);
(2)~the most recent prior reposter (temporal proximity); (3)~the cascade root
(conservative fallback).

Each cascade edge $(s \to t)$ is labelled by cross-referencing
\texttt{followers.csv.gz}: \textbf{Type~0} (follower---organic),
\textbf{Type~1} (non-follower---algorithmic/external), or \textbf{Type~2}
(unknown).  The follower graph is a single snapshot ($\sim$Mar~2024); some edges
may not have existed at cascade time, acknowledged as a limitation.

Cascade size thresholds ($\geq$5 or $\geq$10 reposts) are determined
empirically after preliminary sampling.  Node features comprise: time delta from
root, RoBERTa sentiment label/score \citep{failla2024bluesky}, cascade depth,
and in/out-degree.  Graph-level features include sentiment composition,
follower-edge ratio, event-window label, and structural descriptors (depth,
breadth, branching factor, Wiener index per \citealt{goel2016structural}).

\subsection{BlueCCGL Architecture}

\paragraph{Extension~1: Feed-Aware Augmentation.}
Three domain-specific augmentation strategies replace generic random
node/edge dropping:
\begin{enumerate}[label=\textbf{(\alph*)}]
  \item \emph{Follower-subgraph extraction}---retain only Type~0 edges
        (simulates purely organic spread);
  \item \emph{Temporal prefix sampling}---retain the earliest 60\% of nodes
        (captures early-stage dynamics before amplification); and
  \item \emph{Bridge-node removal}---remove 30\% of non-follower entry points
        (simulates spread without cross-community bridges).
\end{enumerate}
Positive pairs are two augmented views of the same cascade; negatives are
in-batch cascades.

\paragraph{Extension~2: Edge-Type Aware R-GCN Encoder.}
Three relation types (follower, non-follower, unknown) with relation-specific
weight matrices and basis decomposition (2~bases for 3~relations).  Multi-layer
R-GCN (2--3~layers) with ReLU and dropout.

\paragraph{Extension~3: Hierarchical Pooling with Sentiment Enrichment.}
DiffPool-based pooling with pool ratios $[0.5, 0.25]$.  At each hierarchy level,
global mean pooling produces a level-specific embedding concatenated with
graph-level sentiment composition (proportion negative/neutral/positive).  This
captures micro (user), meso (community), and macro (global shape) patterns.

\paragraph{Loss Function.}
InfoNCE contrastive loss following CCGL:
\begin{equation}
\mathcal{L} = -\log \frac{\exp(\mathrm{sim}(z_i, z_i^+) / \tau)}%
              {\sum_{j=1}^{N} \exp(\mathrm{sim}(z_i, z_j) / \tau)}
\end{equation}
where $z_i$ and $z_i^+$ are embeddings of two augmented views of the same
cascade and $\tau$ is the temperature parameter.

\subsection{Archetype Discovery and Evaluation}

\textbf{HDBSCAN} is applied to learned embeddings with minimum cluster size
tuned via grid search (50, 100, 200).  Each discovered archetype is
characterised by its structural profile (depth, breadth, Wiener index),
edge-type composition, sentiment composition, temporal profile, event-window
distribution, and cascade size distribution.

\paragraph{Ablation studies.}
Five ablations isolate the contribution of each extension:
(A1)~replace R-GCN with homogeneous GCN;
(A2)~replace feed-aware augmentations with random node/edge dropping;
(A3)~replace hierarchical pooling with global mean pooling;
(A4)~remove sentiment features;
(A5)~skip contrastive pre-training and cluster raw structural features.

\paragraph{Downstream transfer.}
The pre-trained encoder is fine-tuned on cascade size prediction and cascade
longevity prediction, compared against random initialisation, vanilla CCGL, and
a feature-engineering baseline.

\paragraph{Statistical robustness.}
All results include bootstrap confidence intervals (1000~resamples) and
stability analysis (10$\times$ 80\% subsamples with Adjusted Rand Index).

% =============================================================================
\section{Data Source}
% =============================================================================

The primary data source is the \citet{failla2024bluesky} Zenodo dataset
(DOI:~\href{https://doi.org/10.5281/zenodo.14669616}{10.5281/zenodo.14669616}):
22\,GB compressed, covering 4,099,699~users ($\sim$81\%), 235,567,116~posts, and
$\sim$144M follower edges (Feb~17, 2023 -- Mar~18, 2024).  Key files:
\texttt{user\_posts.tar.gz} (19.5\,GB; text, timestamps, sentiment,
repost metadata), \texttt{interactions.csv.gz} (1.0\,GB; timestamped
edges), \texttt{followers.csv.gz} (491\,MB; follower graph), and
\texttt{feed\_posts.tar.gz} (16.4\,MB; feed-generator posts).

\paragraph{Limitations.}
The follower graph is a single snapshot ($\sim$Mar~2024); some follow
relationships may not have existed during earlier cascades.  User IDs are
pseudonymised integers (no re-identification possible).  Sentiment labels are
available only for English posts; non-English posts receive default neutral.
The English-dominant, left-leaning user base
\citep{quelle2025bluesky} means findings describe crisis dynamics within a
politically homogeneous decentralised community, not global discourse.

% =============================================================================
\section{Expected Outcomes}
% =============================================================================

\paragraph{Outcome~1---A taxonomy of crisis cascade archetypes.}
We expect HDBSCAN clustering to reveal distinct archetypes characterised by
interpretable structural and sentiment properties.  Based on prior literature
\citep{goel2016structural, vosoughi2018spread}, we anticipate at minimum
\emph{broadcast-like} cascades (shallow, wide, high non-follower ratio) and
\emph{viral-like} cascades (deep, narrow, high follower ratio), with potential
additional archetypes (e.g., deliberative, ceasefire/hope, sustained-advocacy).
The key finding is the temporal narrative: how archetype distributions shift
across the six mega-events as the conflict matures.

\paragraph{Outcome~2---Algorithmic vs.\ organic amplification of crisis content.}
By quantifying the follower-edge ratio across archetypes and mega-events, we
provide the first direct measurement of how crisis content spreads through
organic networks versus algorithmically-mediated pathways on a decentralised
platform---a finding structurally impossible to obtain from centralised
platforms.

\paragraph{Outcome~3---Sentiment-structure coupling.}
By embedding sentiment composition into cascade representations, we reveal
whether affective dynamics are coupled to structural propagation patterns.

\paragraph{Outcome~4---Quantitative evidence for each extension.}
Ablation studies (A1--A5) demonstrate the contribution of edge-type encoding,
feed-aware augmentation, hierarchical pooling, and sentiment enrichment.
Negative results are themselves valuable findings.

\paragraph{Outcome~5---Open-source implementation and reusable dataset.}
The project produces: (a)~the BlueCCGL PyTorch Geometric codebase (NLP
filtering, cascade construction, model training, clustering); and (b)~a
processed Bluesky crisis cascade graph dataset with edge-type annotations,
sentiment labels, and event-window assignments.

% =============================================================================
\section{Research Plan}
% =============================================================================

This is a solo Bachelor's thesis project spanning 12~weeks.  The compressed
timeline is enabled by leveraging the pre-existing Zenodo dataset and
pre-computed sentiment labels.
Table~\ref{tab:schedule} details the planned activities.

\begin{longtable}{
  @{}
  >{\centering\arraybackslash}p{0.8cm}
  >{\raggedright\arraybackslash}p{2.6cm}
  >{\raggedright\arraybackslash}p{7.2cm}
  >{\raggedright\arraybackslash}p{4.0cm}
  @{}
}
\caption{Planned weekly schedule for the BlueCCGL project.}
\label{tab:schedule} \\
\toprule
\textbf{Wk} & \textbf{Phase} & \textbf{Activities} & \textbf{Deliverable} \\
\midrule
\endfirsthead
\toprule
\textbf{Wk} & \textbf{Phase} & \textbf{Activities} & \textbf{Deliverable} \\
\midrule
\endhead
\midrule
\multicolumn{4}{r}{\emph{Continued on next page}} \\
\endfoot
\bottomrule
\endlastfoot

1 & Data Acquisition &
  Download Zenodo dataset ($\sim$22\,GB); decompress; load into Parquet/DuckDB;
  preliminary keyword sampling (2--3 of 100 files); validate precision/recall
  on $\sim$200 posts &
  Raw dataset loaded; keyword validation report \\
\addlinespace
2 & NLP Filtering &
  Full-corpus keyword filtering; event-window assignment (E1--E6 + baseline);
  cascade root identification &
  Filtered conflict post index \\
\addlinespace
3 & Cascade Construction &
  Reconstruct cascades; assign edge types via \texttt{followers.csv.gz};
  determine size thresholds &
  Edge-type labelled cascades \\
\addlinespace
4 & Cascade Construction &
  Convert to PyG dataset with node and graph features; detailed EDA &
  PyG dataset; EDA notebook \\
\addlinespace
5 & Model Implementation &
  Base CCGL reproduction: GNN encoder, sub-cascade augmentation, InfoNCE loss;
  Extension~1 (feed-aware augmentations) &
  Verified training loop; augmentation module \\
\addlinespace
6 & Model Implementation &
  Extension~2 (R-GCN encoder); Extension~3 (hierarchical DiffPool + sentiment);
  full BlueCCGL assembly &
  Complete model; training verified \\
\addlinespace
7--8 & Experiments &
  Contrastive pre-training; hyperparameter tuning; HDBSCAN clustering;
  clustering metrics; UMAP visualisations &
  Trained model; archetype clusters \\
\addlinespace
9 & Experiments &
  Ablation studies (A1--A5); downstream transfer evaluation; bootstrap CIs;
  stability analysis &
  Ablation + downstream results \\
\addlinespace
10 & Analysis &
  Archetype interpretation; temporal narrative across mega-events;
  sentiment-structure coupling analysis &
  Archetype characterisation document \\
\addlinespace
11 & Writing &
  Draft all thesis chapters; compile bibliography; create figures and tables &
  Thesis draft \\
\addlinespace
12 & Writing &
  Revisions based on supervisor feedback; final proofreading; prepare
  presentation &
  Final thesis; presentation \\

\end{longtable}

% =============================================================================
\section{Ethical Considerations}
% =============================================================================

\begin{enumerate}[label=\textbf{(\arabic*)}]
  \item \textbf{Data provenance.} The \citet{failla2024bluesky} dataset is
        publicly available under Bluesky's Terms of Service.  All user IDs are
        pseudonymised integers; re-identification is not possible.
  \item \textbf{Conflict sensitivity.} This study analyses the
        \emph{structural dynamics} of information spread, not the content or
        validity of political positions.  Framing is neutral and structural
        throughout.
  \item \textbf{Demographic bias.} All findings are contextualised as
        describing crisis dynamics within Bluesky's specific user base
        (English-dominant, left-leaning early adopters).
  \item \textbf{Responsible disclosure.} Results are reported at the aggregate
        level; no individual cascades are reproduced or identifiable.
\end{enumerate}

% =============================================================================
% Bibliography
% =============================================================================

\newpage
\printbibliography

\end{document}
